\documentclass[12pt, a4paper]{article}
\usepackage[utf8]{inputenc}
\usepackage[T1]{fontenc}
\usepackage{geometry}
\usepackage{enumitem}
\usepackage{titlesec}
\usepackage{hyperref}
\usepackage{xcolor}
\usepackage{booktabs}
\usepackage{array}
\usepackage{fancyhdr}
\usepackage{graphicx}
\usepackage{parskip}

\geometry{
  top=2.5cm,
  bottom=2.5cm,
  left=3cm,
  right=2.5cm
}

% Cores
\definecolor{azulprincipal}{RGB}{70, 130, 180}
\definecolor{azulescuro}{RGB}{25, 25, 112}
\definecolor{cinza}{RGB}{100, 100, 100}

% Cabeçalho e rodapé
\pagestyle{fancy}
\fancyhf{}
\fancyhead[L]{\small\color{cinza} Gerador de Fugas Musicais}
\fancyhead[R]{\small\color{cinza} Desenvolvimento de Software}
\fancyfoot[C]{\thepage}
\renewcommand{\headrulewidth}{0.4pt}

% Formatação de títulos
\titleformat{\section}{\large\bfseries\color{azulescuro}}{}{0em}{\thesection\quad}
\titleformat{\subsection}{\normalsize\bfseries\color{azulprincipal}}{}{0em}{\thesubsection\quad}

\hypersetup{
  colorlinks=true,
  linkcolor=azulescuro,
  urlcolor=azulprincipal
}

% -------------------------------------------------------
\begin{document}

% Capa
\begin{titlepage}
  \centering
  \vspace*{2cm}
  {\Large\bfseries\color{azulescuro} Trabalho de Desenvolvimento de Software\par}
  \vspace{0.5cm}
  {\large\color{cinza} Arthur Rodrigues Padilha, Ester de Lemos Crestani,\\ Isaac de Brito Pereira, Leônidas Lewy, responsavelA\par}
  \vspace{2cm}
  {\Huge\bfseries\color{azulprincipal} Gerador de Fugas Musicais\par}
  \vspace{1cm}
  {\large\bfseries Lista de Requisitos do Sistema\par}
  \vspace{0.5cm}
  {\normalsize 2\textordfeminine{} Fase do Trabalho\par}
  \vfill
  {\normalsize Porto Alegre, Maio de 2026}
\end{titlepage}

% Sumário
\tableofcontents
\newpage

% -------------------------------------------------------
\section{Visão Geral do Sistema}

O \textbf{Gerador de Fugas Musicais} é um software que converte texto em saída sonora, mapeando cada caractere para notas, pausas e alterações de configuração, focando agora na geração de música polifônica em formato de "fuga" (estilo Johann Sebastian Bach). Cada linha do texto representa uma voz independente tocada simultaneamente. O sistema suporta leitura/gravação de arquivos TXT e exportação para arquivos MIDI.

\subsection{Tecnologias Utilizadas}

\begin{itemize}[leftmargin=2cm]
  \item \textbf{Linguagem:} Python 3.12
  \item \textbf{Interface Gráfica:} CustomTkinter (substituiu Tkinter padrão para visual moderno)
  \item \textbf{Gerenciador de dependências:} UV
  \item \textbf{Versionamento:} Git e GitHub
  \item \textbf{Framework de testes:} PyTest
  \item \textbf{Áudio em tempo real:} pygame.midi
  \item \textbf{Geração de arquivos MIDI:} mido
\end{itemize}

% -------------------------------------------------------
\section{Requisitos Funcionais}

Os requisitos funcionais descrevem as funcionalidades que o sistema deve oferecer ao usuário.

\subsection{RF01 — Entrada de Texto pelo Usuário}

O sistema deve disponibilizar uma interface com uma caixa de texto onde o usuário possa digitar livremente sua sequência de caracteres ou carregá-la de um arquivo.

\begin{itemize}[leftmargin=2cm]
  \item A entrada pode ser feita via digitação ou carregamento de arquivo TXT.
  \item O campo deve ser editável.
  \item Deve ser possível salvar o texto modificado de volta em um arquivo TXT.
  \item Cada linha no texto representa uma \textbf{voz musical independente}.
\end{itemize}

\subsection{RF02 — Interpretação da Sequência em Fuga}

O sistema deve processar a entrada onde cada linha (voz) tem configurações iniciais cíclicas de oitava, volume e instrumento (Voz 0: Oitava 6, Vol 100, Piano; Voz 1: Oitava 5, Vol 80, Órgão, etc.).

\begin{center}
\begin{tabular}{>{\bfseries}p{3.5cm} p{9cm}}
  \toprule
  Caractere(s) & Comportamento \\
  \midrule
  A a G (maiúsc.) & Notas musicais Lá a Sol \\
  H & Si bemol \\
  Mb & Mi bemol \\
  a--h (minúsc.) & Silêncio/Pausa \\
  ! & Altera instrumento para Harmonica (MIDI 22) \\
  ; & Altera instrumento para Tubular Bells (MIDI 14) \\
  , & Altera instrumento para Church Organ (MIDI 19) \\
  O, I, U & Altera instrumento para Gaita de Foles/Bagpipe (MIDI 109) \\
  Espaço & Dobra o volume da voz atual \\
  ? & Aumenta a oitava da voz atual \\
  V & Diminui a oitava da voz atual \\
  {>} & Aumenta o BPM global em 10 \\
  {<} & Diminui o BPM global em 10 \\
  {[}n{]} no início da linha & Atraso de \textit{n} beats antes da voz começar \\
  Outras letras / símbolos & Repetem a nota anterior; se não houver nota anterior, geram Silêncio \\
  \bottomrule
\end{tabular}
\end{center}

\subsection{RF03 — Geração e Reprodução Polifônica}

O sistema deve reproduzir os tokens de todas as vozes em tempo real e de forma sincronizada (polifonia). A classe \texttt{ReproducaoAudio} recebe notas, volume, oitava e o canal MIDI correspondente.

\subsection{RF04 — Exportação de Arquivo MIDI}

Além da saída sonora, o sistema deve salvar toda a música gerada em um arquivo \texttt{.mid} padrão ao final da reprodução, na pasta \texttt{saida/}.

\subsection{RF05 — Controles de Interface Gráfica (5 Telas)}

O sistema possui um fluxo de navegação entre telas:
\begin{itemize}[leftmargin=2cm]
  \item \textbf{Tela Inicial:} Botão ``Começar'' e modal ``Informações'' com a legenda de mapeamento.
  \item \textbf{Tela de Entrada:} Carregamento de TXT e digitação da sequência, com legenda detalhada visível.
  \item \textbf{Tela de Configurações:} Configuração por voz (instrumento, oitava, BPM, volume) com controles interativos.
  \item \textbf{Tela de Reprodução:} Exibe o estado em tempo real de cada voz (nota atual, oitava, instrumento, volume, status).
  \item \textbf{Tela Final:} Informa o caminho do arquivo \texttt{.mid} salvo e botão para nova sequência.
\end{itemize}


% -------------------------------------------------------
\section{Requisitos Não Funcionais}

\subsection{RNF01 — Linguagem e Plataforma}
O sistema deve ser implementado em Python, utilizando UV para dependências e Git/GitHub para versionamento.

\subsection{RNF02 — Modularidade e Testabilidade (SOLID)}
As classes devem possuir responsabilidade única, com testes unitários (PyTest) desenvolvidos simultaneamente à criação das classes (ex: \texttt{teste\_maestro.py}).

\subsection{RNF03 — Compatibilidade}
Deve utilizar \texttt{pygame.midi} para o áudio em tempo real e \texttt{mido} para a geração dos arquivos `.mid`.

% -------------------------------------------------------
\section{Modelagem das Classes Principais}

\subsection{Classe \texttt{Voz}}
Encapsula o estado independente de uma linha de texto (uma voz musical). Controla sua própria oitava, volume, canal e instrumento.
\begin{itemize}[leftmargin=2cm]
  \item \textbf{Atributos:} \texttt{id\_voz: int}, \texttt{linha\_texto: str}, \texttt{oitava\_atual: int}, \texttt{volume\_atual: int}, \texttt{instrumento\_atual: int}, \texttt{canal: int}, \texttt{atraso\_inicial: int}, \texttt{nota\_anterior: str}.
  \item \textbf{Métodos Principais:}
    \begin{itemize}
      \item \texttt{\_parse\_inicio()} \textbf{->} Método interno que lê o início da linha reconhecendo, em sequência: caractere de instrumento (\texttt{!}, \texttt{;}, \texttt{,}), atraso \texttt{[n]}, e posiciona o leitor no primeiro token de nota.
      \item \texttt{ler\_proximo\_token()} \textbf{->} Retorna um \texttt{dict} com tipo de evento (NOTA, PAUSA, INSTRUMENTO, BPM\_UP, etc.) e o caractere lido.
      \item \texttt{has\_next()} \textbf{->} Retorna \texttt{bool} indicando se ainda há caracteres a processar.
      \item \texttt{reiniciar\_estado()} \textbf{->} Redefine a leitura para o início.
    \end{itemize}
  \item \textbf{Entrada:} Inteiro (ID da voz, define parâmetros cíclicos) e string (linha de texto).
  \item \textbf{Saída:} Tokens sequenciais (dicionários) para o Maestro executar.
\end{itemize}

\subsection{Classe \texttt{Interpretador} (Partitura)}
Responsável por receber o texto completo da composição, dividi-lo em vozes e gerenciar essas instâncias.
\begin{itemize}[leftmargin=2cm]
  \item \textbf{Atributos:} \texttt{vozes: List[Voz]}.
  \item \textbf{Métodos Principais:}
    \begin{itemize}
      \item \texttt{\_inicializar\_vozes(texto: str)} \textbf{->} Instancia as vozes quebrando o texto por quebra de linhas.
      \item \texttt{get\_vozes()} \textbf{->} Retorna a lista de objetos \texttt{Voz}.
    \end{itemize}
  \item \textbf{Entrada:} Uma única string (\texttt{str}) com o texto cru completo digitado pelo usuário.
  \item \textbf{Saída:} Uma coleção estruturada de objetos \texttt{Voz} configurados.
\end{itemize}

\subsection{Classe \texttt{Maestro}}
Orquestrador central do sistema. Roda o loop musical processando todas as vozes simultaneamente (polifonia) e controlando o tempo.
\begin{itemize}[leftmargin=2cm]
  \item \textbf{Atributos:} \texttt{reproducao: ReproducaoAudio}, \texttt{gerador\_midi: GeradorMIDI}, \texttt{bpm: int}, \texttt{partitura: Interpretador}.
  \item \textbf{Métodos Principais:}
    \begin{itemize}
      \item \texttt{preparar(texto: str)} \textbf{->} Inicializa a Partitura e prepara os geradores.
      \item \texttt{tocar()} \textbf{->} Inicia a \textit{thread} de reprodução e envia configurações iniciais aos canais.
      \item \texttt{\_loop\_musical()} \textbf{->} Avança beat por beat, extraindo um token de cada voz e delegando o comando de tocar ou pausar. Controla a thread e os sleeps.
    \end{itemize}
  \item \textbf{Entrada:} Recebe comandos de interface (texto bruto na preparação) e os tokens consumidos da Partitura.
  \item \textbf{Saída:} Chamadas de funções acopladas para enviar dados estruturados (nota, canal, volume) para as classes de Áudio e MIDI.
\end{itemize}

\subsection{Classe \texttt{ReproducaoAudio}}
Responsável por interagir diretamente com a API do sistema operacional (\texttt{pygame.midi}) para produzir sons multitimbrais em tempo real.
\begin{itemize}[leftmargin=2cm]
  \item \textbf{Atributos:} \texttt{\_output: pygame.midi.Output}.
  \item \textbf{Métodos Principais:}
    \begin{itemize}
      \item \texttt{reproduzirNota(nota: str, oitava: int, volume: int, canal: int)} \textbf{->} Toca a nota no hardware.
      \item \texttt{silenciarNota(nota: str, oitava: int, canal: int)} \textbf{->} Para o som emitindo Note Off.
      \item \texttt{setInstrumento(id\_midi: int, canal: int)} \textbf{->} Altera o timbre do canal.
    \end{itemize}
  \item \textbf{Entrada:} Instruções atômicas enviadas pelo Maestro, contendo a string da nota exata, seu volume dinâmico, sua oitava e o canal alvo.
  \item \textbf{Saída:} Sinais de áudio para os alto-falantes por meio da biblioteca MIDI subjacente.
\end{itemize}

\subsection{Classes de Suporte}
\begin{itemize}[leftmargin=2cm]
  \item \textbf{\texttt{GeradorMIDI}:} Transforma os eventos despachados pelo Maestro em mensagens do formato \texttt{mido} e salva o arquivo final \texttt{.mid}. Recebe notas/tempos exatos e entrega um arquivo de mídia compilado.
  \item \textbf{\texttt{GerenciadorArquivo}:} Fornece métodos estáticos para \texttt{ler\_arquivo\_texto(caminho)} e \texttt{salvar\_arquivo\_texto(caminho, conteudo)}. Recebe strings e caminhos e atua diretamente no sistema de arquivos.
\end{itemize}

% -------------------------------------------------------
\section{Divisão de Responsabilidades}

\begin{center}
\begin{tabular}{>{\bfseries}p{4cm} p{9cm}}
  \toprule
  Tarefa & Responsáveis \\
  \midrule
  Lista de requisitos & Ester Crestani (Atualização), Isaac (Revisão) \\
  Lista de classes & Ester Crestani (Atualização), Isaac (Revisão) \\
  Interfaces & Ester Crestani (Atualização e Formatação visual) \\
  Front-end & Ester Crestani \\
  Main (back-end) & responsavelA \\
  Classe \texttt{Voz} & Ester Crestani \\
  Classe \texttt{Interpretador} & responsavelA \\
  Classe \texttt{Maestro} & Arthur \\
  Classe \texttt{ReproducaoAudio} & Isaac de Brito Pereira \\
  Testes & Desenvolvidos pelos responsáveis de cada classe \\
  \bottomrule
\end{tabular}
\end{center}

\end{document}
